\subsection{Householder Reflections}

Throughout this subsection, $\mathbb{F}\in\{\mathbb{R},\mathbb{C}\}$.

\begin{definition} \label{def:householder-matrix} \hlt{Householder Matrix}\\
Let $n\geq1$ and $\mathbf{w}\in\mathbb{F}^n\setminus\{\mathbf{0}\}$, and set $\mathbf{u}\defeq\frac{\mathbf{w}}{\norm{\mathbf{w}}_2}$. The \hlt{Householder matrix} associated with $\mathbf{w}$ is $U_{\mathbf{w}}\defeq I_n-2P_{\mathbf{u}}=I_n-2\frac{\mathbf{w}\mathbf{w}^*}{\norm{\mathbf{w}}_2^2}$.\\
Equivalently, $(\forall\,\mathbf{x}\in\mathbb{F}^n)\,U_{\mathbf{w}}\mathbf{x}=\mathbf{x}-2\innerprod[\mathbb{F}^n]{\mathbf{x},\mathbf{u}}\mathbf{u}$.
\end{definition}

\begin{theorem} \label{thm:householder-matrix-properties} \hlt{Householder Matrices Are Unitary Reflections}\\
Let $n\geq1$ and let $\mathbf{w}\in\mathbb{F}^n\setminus\{\mathbf{0}\}$. Then:
\begin{enumerate}[label=(\roman*)]
\setlength{\itemsep}{0pt}
\item $U_{\mathbf{w}}^*=U_{\mathbf{w}}$;
\item $U_{\mathbf{w}}^2=I_n$;
\item $U_{\mathbf{w}}^*U_{\mathbf{w}}=I_n$, so $U_{\mathbf{w}}$ is unitary;
\item $(\forall\,\mathbf{x}\in\mathbf{w}^{\perp})\,U_{\mathbf{w}}\mathbf{x}=\mathbf{x}$;
\item $(\forall\,\mathbf{x}\in\operatorname{span}(\mathbf{w}))\,U_{\mathbf{w}}\mathbf{x}=-\mathbf{x}$;
\item $\mathbf{w}\in\mathbb{R}^n\Longrightarrow U_{\mathbf{w}}$ is symmetric and orthogonal over $\mathbb{R}$.
\end{enumerate}
\begin{proof}
Set $\mathbf{u}\defeq\frac{\mathbf{w}}{\norm{\mathbf{w}}_2}$ and $P\defeq P_{\mathbf{u}}$.
\begin{enumerate}[label=(\roman*)]
\setlength{\itemsep}{0pt}
\item The identity $P^*=P$ holds by Proposition \ref{prop:rank-one-projection-matrix-properties}(iii). Hence $U_{\mathbf{w}}^*=I_n-2P=U_{\mathbf{w}}$.
\item The identity $P^2=P$ holds by Proposition \ref{prop:rank-one-projection-matrix-properties}(iv). Hence $U_{\mathbf{w}}^2=(I_n-2P)^2=I_n-4P+4P^2=I_n$.
\item The identity $U_{\mathbf{w}}^*U_{\mathbf{w}}=U_{\mathbf{w}}^2=I_n$ holds by (i)--(ii), so $U_{\mathbf{w}}$ is unitary.
\item Fix $\mathbf{x}\in\mathbf{w}^{\perp}=\operatorname{span}(\mathbf{u})^{\perp}$. Then $P\mathbf{x}=\mathbf{0}$ by Proposition \ref{prop:rank-one-projection-matrix-properties}(ii), so $U_{\mathbf{w}}\mathbf{x}=\mathbf{x}$.
\item Fix $\mathbf{x}\in\operatorname{span}(\mathbf{w})=\operatorname{span}(\mathbf{u})$, and choose $a\in\mathbb{F}$ with $\mathbf{x}=a\mathbf{u}$. Then $P\mathbf{x}=\innerprod{a\mathbf{u},\mathbf{u}}\mathbf{u}=a\mathbf{u}=\mathbf{x}$ by $\norm{\mathbf{u}}_2=1$, so $U_{\mathbf{w}}\mathbf{x}=-\mathbf{x}$.
\item Suppose $\mathbf{w}\in\mathbb{R}^n$. Then $\mathbf{u}$, $P$, and $U_{\mathbf{w}}$ have real entries.\\
The identities $U_{\mathbf{w}}^t=U_{\mathbf{w}}$ and $U_{\mathbf{w}}^tU_{\mathbf{w}}=I_n$ hold by (i) and (iii), respectively.
\end{enumerate}
\hlt{Structural Commentary}: Self-adjoint idempotence proves unitarity; the range--kernel split identifies the two eigenspaces.\\
\hlt{Key Insight}: $I_n-2P_{\mathbf{u}}$ fixes $\ker(P_{\mathbf{u}})$ and negates $\operatorname{range}(P_{\mathbf{u}})$.
\end{proof}
\end{theorem}

Figure \ref{fig:householder-reflection} illustrates the two eigenspace actions in Theorem \ref{thm:householder-matrix-properties}(iv)--(v).
\begin{illustrationfigure}[H]
  \begin{tikzpicture}[
  >=Stealth,
  every node/.style={font=\small},
  plane/.style={fill=gray!18, draw=gray!45, line width=0.7pt},
  normal/.style={->, black, line width=0.9pt},
  vec/.style={->, blue!75!black, line width=0.95pt},
  refl/.style={->, teal!75!black, line width=0.95pt},
  guide/.style={gray!60, dashed, line width=0.65pt}
]
  \coordinate (O) at (0,0);
  \coordinate (A) at (-2.2,-0.45);
  \coordinate (B) at (2.6,-0.45);
  \coordinate (C) at (1.75,0.55);
  \coordinate (D) at (-3.05,0.55);
  \coordinate (X) at (1.45,2.0);
  \coordinate (PX) at (0.72,0.05);
  \coordinate (RX) at (0.0,-1.9);

  \draw[plane] (A) -- (B) -- (C) -- (D) -- cycle;
  \node at (2.35,0.35) {$\mathbf{w}^{\perp}$};

  \draw[normal] (O) -- (0,2.2) node[above] {$\mathbf{w}$};
  \draw[normal] (O) -- (0,-1.05) node[left] {$-\mathbf{w}$};

  \draw[vec] (O) -- (X) node[above right] {$\mathbf{x}$};
  \draw[refl] (O) -- (RX) node[below] {$U_{\mathbf{w}}\mathbf{x}$};
  \draw[guide] (X) -- (PX);
  \draw[guide] (RX) -- (PX);
  \draw[gray!70] ($(PX)+(-0.14,0.08)$) -- ($(PX)+(0.02,0.2)$) -- ($(PX)+(0.16,0.08)$);
\end{tikzpicture}

  \caption{Householder reflection: $U_{\mathbf{w}}$ preserves the $\mathbf{w}^{\perp}$ component of $\mathbf{x}$ and negates its $\operatorname{span}(\mathbf{w})$ component.}
  \label{fig:householder-reflection}
\end{illustrationfigure}

\begin{theorem} \label{thm:householder-maps-equal-norm-vectors} \hlt{Householder Map Between Equal-Norm Vectors}\\
Let $n\geq1$ and let $\mathbf{x},\mathbf{y}\in\mathbb{F}^n$ satisfy $\norm{\mathbf{x}}_2=\norm{\mathbf{y}}_2\neq0$. Set $z\defeq\innerprod[\mathbb{F}^n]{\mathbf{x},\mathbf{y}}$.\\
Choose $\sigma\in\mathbb{F}$ with $\sigma=1$ for $z=0$ and $\sigma=-\frac{\overline{z}}{\abs{z}}$ for $z\neq0$. Set $\mathbf{w}\defeq\mathbf{y}-\sigma\mathbf{x}$. Then:
\begin{enumerate}[label=(\roman*)]
\setlength{\itemsep}{0pt}
\item $\abs{\sigma}=1$;
\item $\mathbf{w}\neq\mathbf{0}$;
\item $\sigma U_{\mathbf{w}}\mathbf{x}=\mathbf{y}$.
\end{enumerate}
\begin{proof}
Set $r\defeq\norm{\mathbf{x}}_2=\norm{\mathbf{y}}_2$.
\begin{enumerate}[label=(\roman*)]
\setlength{\itemsep}{0pt}
\item The identity $\abs{\sigma}=1$ holds by the two defining cases.
\item The identity $\sigma z=-\abs{z}$ holds by the two defining cases.\\
Thus $\norm{\mathbf{w}}_2^2=2r^2-2\operatorname{Re}(\sigma z)=2(r^2+\abs{z})>0$ by (i) and conjugate symmetry. Hence $\mathbf{w}\neq\mathbf{0}$ by positive definiteness.
\item By $\mathbf{y}=\mathbf{w}+\sigma\mathbf{x}$ and Definition \ref{def:householder-matrix}, $\sigma U_{\mathbf{w}}\mathbf{x}=\mathbf{y}-\mathbf{w}-2\sigma\frac{\innerprod{\mathbf{x},\mathbf{w}}}{\norm{\mathbf{w}}_2^2}\mathbf{w}$.\\
The identity $2\sigma\innerprod{\mathbf{x},\mathbf{w}}=2\sigma z-2r^2=-\norm{\mathbf{w}}_2^2$ holds by the definition of $\mathbf{w}$ and (i)--(ii). Therefore $\sigma U_{\mathbf{w}}\mathbf{x}=\mathbf{y}$.
\end{enumerate}
\hlt{Structural Commentary}: The phase makes $\sigma z$ nonpositive real; the norm identity then cancels the reflected component.\\
\hlt{Key Insight}: $2\sigma\innerprod{\mathbf{x},\mathbf{w}}=-\norm{\mathbf{w}}_2^2$ annihilates the residual $\mathbf{w}$-coefficient.
\end{proof}
\end{theorem}

\begin{corollary} \label{cor:householder-unitary-annihilator} \hlt{Householder Unitary Annihilator}\\
Let $n\geq1$ and $\mathbf{x}=[x_1\ \cdots\ x_n]^t\in\mathbb{F}^n\setminus\{\mathbf{0}\}$. Choose $\sigma\in\mathbb{F}$ with $\sigma=1$ for $x_1=0$ and $\sigma=-\frac{\overline{x_1}}{\abs{x_1}}$ for $x_1\neq0$.\\
Set $\mathbf{w}\defeq\norm{\mathbf{x}}_2\mathbf{e}_1-\sigma\mathbf{x}$. Then:
\begin{enumerate}[label=(\roman*)]
\setlength{\itemsep}{0pt}
\item $\mathbf{w}\neq\mathbf{0}$;
\item $\sigma U_{\mathbf{w}}\mathbf{x}=\norm{\mathbf{x}}_2\mathbf{e}_1$;
\item $\norm{\mathbf{x}}_2=1\Longrightarrow\overline{\sigma}U_{\mathbf{w}}$ is unitary with $(\overline{\sigma}U_{\mathbf{w}})\mathbf{e}_1=\mathbf{x}$.
\end{enumerate}
\begin{proof}
Set $\mathbf{y}\defeq\norm{\mathbf{x}}_2\mathbf{e}_1$. Then $\norm{\mathbf{y}}_2=\norm{\mathbf{x}}_2$ and $\innerprod{\mathbf{x},\mathbf{y}}=x_1\norm{\mathbf{x}}_2$.\\
Since $\norm{\mathbf{x}}_2>0$, the defining cases for $\sigma$ coincide with those determined by $\innerprod{\mathbf{x},\mathbf{y}}$.
\begin{enumerate}[label=(\roman*)]
\setlength{\itemsep}{0pt}
\item The relation $\mathbf{w}\neq\mathbf{0}$ holds by Theorem \ref{thm:householder-maps-equal-norm-vectors}(ii) with this $\mathbf{y}$.
\item The identity $\sigma U_{\mathbf{w}}\mathbf{x}=\norm{\mathbf{x}}_2\mathbf{e}_1$ holds by Theorem \ref{thm:householder-maps-equal-norm-vectors}(iii) with this $\mathbf{y}$.
\item Suppose $\norm{\mathbf{x}}_2=1$ and set $A\defeq\sigma U_{\mathbf{w}}$. The identity $A^*A=\abs{\sigma}^2U_{\mathbf{w}}^*U_{\mathbf{w}}=I_n$ holds by Theorem \ref{thm:householder-maps-equal-norm-vectors}(i) and Theorem \ref{thm:householder-matrix-properties}(iii).\\
Thus $A^*$ is unitary by Theorem \ref{thm:unitary-matrix-equivalent-conditions}(i),(iv) with $U=A$. By (ii), $A\mathbf{x}=\mathbf{e}_1$, so $A^*\mathbf{e}_1=A^*A\mathbf{x}=\mathbf{x}$.\\
The identity $A^*=\overline{\sigma}U_{\mathbf{w}}$ holds by Theorem \ref{thm:householder-matrix-properties}(i).
\end{enumerate}
\hlt{Structural Commentary}: Equal-norm reflection constructs $A\mathbf{x}=\mathbf{e}_1$; unitarity transfers this identity to the first column of $A^*$.\\
\hlt{Key Insight}: The adjoint reverses the annihilation map and installs $\mathbf{x}$ as a column.
\end{proof}
\end{corollary}
